\begin{afterword}
\summaryhead

The post recommends George Orwell's ``Politics and the English Language'' as required reading for rationalists. Orwell's subject was how evil hides in muddy language: a regime that imprisons people without trial says ``Unreliable elements were subjected to an alternative justice process.'' The author illustrates with the passive voice, which removes the actor, and with the passive in scientific papers, where ``The subjects were administered Progenitorivox'' sounds authoritative and hides who handed a student a bottle of pills.

To the objection that a thoughtful reader knows the scientist is there, the author answers with a writer's view. Readers do not stop to think; a statement's public effect is the first impression it makes, and whatever the audience thinks you said is what you said. A rationalist must therefore notice the experience words create, beyond their literal meaning: ``Meaning does not excuse impact.'' The post ends with two passages from Orwell, on speakers who repeat stock phrases and on letting the meaning choose the word.

\responsehead

The post's advice has a sound core, and it credits its source. Attend to what your words make a reader see, and do not excuse a misleading phrase because a careful reading can rescue it. The post names Orwell, quotes him accurately, and its example of euphemism is Orwell's, adapted.

Its account of the passive voice is looser than its advice. ``Passive voice removes the actor'' is true only of passives without a ``by'' phrase, and an active sentence can hide an actor too. Orwell's own example shows that voice is not the problem. His plain statement of the facts, ``People are imprisoned for years without trial,'' is itself a passive with no actor, and he offers it as the version that calls up mental pictures; the euphemism is the noun phrase ``elimination of unreliable elements.'' The post's remark about ``static noun phrases'' comes closer to this than its rule about voice. For journal articles, it explains the passive by its authoritative sound and does not mention that some scientific bodies asked for it.

The main argument changes whose words are being judged. ``Whatever the audience thinks you said is what you said'' is a rule for one's own words, as the author's reply to a commenter put it: misreading is ``always the author's fault, at least if you happen to be an author.'' The post then uses the rule to judge other people's statements, by the impression they make rather than by what they could reasonably mean. That can be a fair standard for public speech. But the post's evidence for an impression is the author's own reaction, reported as what every reader feels (``the only experience''), and its one case drawn from a real speaker is the Summit speaker from ``Applause Lights,'' whose remark the notes on that post discuss.

The post also breaks its own rule. Authors, it says, ``are trained not to give themselves the benefit of the doubt,'' which does not say who trains them, the construction it had just demonstrated with ``Writers are told.'' After a reader asked ``Trained by whom?'', the author's next comment was ``Yes, well, I never said I was a True Writing Master either.''

In short: sound advice about the effect of words, credited to Orwell, with an account of the passive voice that Orwell's own example does not support, and a standard for judging other people's words whose only measure is the author's impression.
\end{afterword}
